\documentclass{beamer}

\input{MathMacros}

\usepackage{beamerthemesplit}
\usepackage{enumerate}
\usepackage{amssymb}
\usepackage[all,cmtip]{xy}
\usepackage[mathscr]{eucal}
\usepackage[natbib=true,style=authoryear,backend=bibtex,useprefix=true]{biblatex}
\addbibresource{reading.bib}
\usepackage{graphicx,color}

\title{Failure Modes of LLM Systems: Why They Break Even When They “Work”}
\author{Reuben Brasher}
\date{\today}

\begin{document}

\frame{\titlepage}

\section[Outline]{}
\frame{\tableofcontents}

\section{Lecture 5}

\frame{
\frametitle{Convincing Failure}
The most dangerous failures are those that appear correct.
\begin{itemize}
\item Outputs are fluent and well-structured.
\item Errors are embedded within plausible explanations.
\item Confidence is simulated rather than earned.
\item Users tend to trust coherent language.
\item Failures require scrutiny to detect.
\item The system fails in ways that resemble success.
\end{itemize}
}

\frame{
\frametitle{Bridge from Phase 2}
Structured workflows improve reliability but do not eliminate failure.
\begin{itemize}
\item Scaffolding improves reliability but not correctness.
\item Transformation pipelines reduce risk but do not remove it.
\item Role separation limits but does not stop error propagation.
\item Systems remain probabilistic.
\item Failures persist under complexity.
\item Reliability remains conditional.
\end{itemize}
}

\frame{
\frametitle{Taxonomy of Failure}
Failures can be categorized into distinct types.
\begin{itemize}
\item Epistemic failures concern truth.
\item Structural failures concern consistency.
\item Process failures concern interaction.
\item Each type requires different mitigation.
\item Failures often overlap categories.
\item Taxonomy enables diagnosis.
\end{itemize}
}

\frame{
\frametitle{Epistemic Failures}
Epistemic failures produce plausible but incorrect outputs.
\begin{itemize}
\item Statements are not grounded in data.
\item Outputs appear analytical.
\item Confidence is unrelated to correctness.
\item Errors require verification to detect.
\item Coherence masks inaccuracy.
\item Truth is not internally validated.
\end{itemize}
}

\frame{
\frametitle{Example: Fabricated Analysis}
Models generate convincing but incorrect analytical conclusions.
\begin{itemize}
\item Trends inferred from insufficient data.
\item Quantitative claims invented.
\item Statistical reasoning simulated.
\item Outputs resemble valid analysis.
\item No numerical verification occurs.
\item External validation is required.
\end{itemize}
}

\frame{
\frametitle{Simulation vs Analysis}
The model simulates reasoning rather than performing it.
\begin{itemize}
\item Pattern matching replaces computation.
\item Language structure substitutes logic.
\item Outputs mimic analytical formats.
\item No execution of statistical methods.
\item Results reflect training patterns.
\item Plausibility replaces correctness.
\end{itemize}
}

\frame{
\frametitle{Implications of Epistemic Failure}
Epistemic failures are dangerous in analytical contexts.
\begin{itemize}
\item Research conclusions may be invalid.
\item Decisions may be misled.
\item Incorrect insights appear authoritative.
\item Novices are especially vulnerable.
\item Errors scale with trust.
\item Verification becomes mandatory.
\end{itemize}
}

\frame{
\frametitle{Fabrication and Hallucination}
Models generate content that appears valid but is invented.
\begin{itemize}
\item Citations may not exist.
\item References may be misattributed.
\item Test outputs may be fabricated.
\item Explanations lack grounding.
\item Content appears complete.
\item Uncertainty is not indicated.
\end{itemize}
}

\frame{
\frametitle{Nature of Hallucination}
Hallucinations follow predictable patterns.
\begin{itemize}
\item Occur when information is missing.
\item Gaps are filled with patterns.
\item Outputs reflect likelihood, not truth.
\item Fluency increases credibility.
\item Errors are systematic.
\item Detection requires validation.
\end{itemize}
}

\frame{
\frametitle{Structural Failures}
Structural failures arise from inconsistency across components.
\begin{itemize}
\item Schema and code may misalign.
\item Representations may diverge.
\item Shared structure may be lost.
\item Errors propagate through stages.
\item Systems become incoherent.
\item Correctness breaks across layers.
\end{itemize}
}

\frame{
\frametitle{Example: Representation Drift}
Complex systems reveal structural instability.
\begin{itemize}
\item JSON differs from implementation.
\item Code violates defined structure.
\item Tests fail from mismatches.
\item Changes do not propagate.
\item Systems become brittle.
\item Alignment is lost over time.
\end{itemize}
}

\frame{
\frametitle{Lack of Internal Consistency}
Models do not enforce consistency across outputs.
\begin{itemize}
\item Each output is independent.
\item No persistent internal model exists.
\item Prior structure is not preserved.
\item Alignment must be enforced externally.
\item Consistency degrades without constraints.
\item Structure is not automatic.
\end{itemize}
}

\frame{
\frametitle{Process Failures}
Process failures arise from the interaction itself.
\begin{itemize}
\item Conversations become unstable.
\item Outputs drift from intent.
\item Irrelevant content appears.
\item Style degrades over time.
\item Context becomes noisy.
\item Workflow loses coherence.
\end{itemize}
}

\frame{
\frametitle{Example: Context Drift}
Extended interaction introduces instability.
\begin{itemize}
\item Instructions are forgotten.
\item Irrelevant artifacts appear.
\item Responses lose focus.
\item Format becomes inconsistent.
\item Errors influence future outputs.
\item Conversation diverges.
\end{itemize}
}

\frame{
\frametitle{Verbosity and Degradation}
Output quality changes over time without control.
\begin{itemize}
\item Responses become longer.
\item Precision decreases.
\item Explanations expand unnecessarily.
\item Formatting becomes inconsistent.
\item Signal-to-noise decreases.
\item User must constrain output.
\end{itemize}
}

\frame{
\frametitle{Interaction Instability}
The interaction itself introduces error.
\begin{itemize}
\item Responses condition future responses.
\item Errors accumulate.
\item Context mixes correct and incorrect data.
\item Model cannot distinguish them.
\item Feedback loops reinforce mistakes.
\item Stability must be managed.
\end{itemize}
}

\frame{
\frametitle{Silent Degradation}
Errors accumulate without immediate detection.
\begin{itemize}
\item Early inaccuracies go unnoticed.
\item Model reuses incorrect outputs.
\item Errors compound.
\item System appears consistent.
\item Detection becomes difficult.
\item Failures emerge late.
\end{itemize}
}

\frame{
\frametitle{User Silence as Reinforcement}
Lack of correction reinforces errors.
\begin{itemize}
\item Incorrect outputs are accepted.
\item Model adapts to perceived approval.
\item Errors persist in future outputs.
\item Mistakes become entrenched.
\item Oversight is required.
\item Correction must be explicit.
\end{itemize}
}

\frame{
\frametitle{Why Failures Occur}
Failures derive from fundamental system limitations.
\begin{itemize}
\item No internal model of truth exists.
\item No persistent state is maintained.
\item No verification mechanism is present.
\item Outputs are generated token-by-token.
\item Correctness is not evaluated.
\item Consistency is not enforced.
\end{itemize}
}

\frame{
\frametitle{Core Conclusion}
Reliability must be imposed rather than assumed.
\begin{itemize}
\item Structure reduces error.
\item Verification must be external.
\item Constraints must be explicit.
\item Oversight must be continuous.
\item Systems must anticipate failure.
\item Trust requires validation.
\end{itemize}
}

\frame{
\frametitle{Practical Implications}
Effective use requires disciplined control.
\begin{itemize}
\item Verify outputs against sources.
\item Use structured inputs.
\item Restart conversations when needed.
\item Limit context to relevance.
\item Treat outputs as drafts.
\item Do not trust fluency.
\end{itemize}
}

\frame{
\frametitle{Operational Rules}
A set of rules mitigates common failures.
\begin{itemize}
\item Ground work in real data.
\item Constrain outputs explicitly.
\item Separate roles and tasks.
\item Validate at each stage.
\item Detect errors early.
\item Maintain workflow control.
\end{itemize}
}

\frame{
\frametitle{Synthesis}
Failures are predictable and systematic.
\begin{itemize}
\item Epistemic failures distort truth.
\item Structural failures break systems.
\item Process failures destabilize workflows.
\item All are inherent to the system.
\item None are eliminated by scale.
\item All require external control.
\end{itemize}
}

\frame{
\frametitle{Transition to Lecture 6}
Failures extend beyond individual use into larger systems.
\begin{itemize}
\item Writing errors become misinformation.
\item Code errors become system failures.
\item Analysis errors become decisions.
\item Risks increase at scale.
\item Systems amplify failure impact.
\item Next topic is security and governance.
\end{itemize}
}

\begin{frame}[t,allowframebreaks]
\frametitle{References}
\printbibliography
\end{frame}

\end{document}